\documentclass[11pt]{article}

\usepackage[a4paper,margin=1in]{geometry}
\usepackage{microtype}
\usepackage{amsmath,amssymb,amsthm}
\usepackage{booktabs}
\usepackage{enumitem}

\newcommand{\R}{\mathbb{R}}
\newcommand{\sg}{\operatorname{sg}}
\newcommand{\norm}[1]{\left\lVert #1 \right\rVert}
\newcommand{\Opi}{\mathrm{O}_\pi}
\newcommand{\Keff}{\mathrm{K}_{\mathrm{eff}}}

\title{Semantic IDs for Generative Recommendation:\\
RQ-VAE, DRQ, and a Controlled Benchmark on TencentGR-1M\\
\large Preface}
\author{rq-vae-msbd-5015}
\date{November 2026}

\begin{document}
\maketitle

\section{In plain terms}

A generative recommender does not look products up in a table. It \emph{writes
out} the identity of the product it wants to show, one token at a time, the way
language is generated. To make that possible, every product in the catalog is
first given a short code built from its content (text, image, or in our case a
precomputed embedding). The codes are designed so that similar products get
similar codes: the first number is a broad group label, the next numbers refine
it. This document studies the two main ways those codes are manufactured today.

Both methods aim at the same business property: every product should end up with
a \emph{private} code. If two products share one code, the model cannot say
which of the two it means, and the pair becomes one coin flip at serving time.
With three code positions and 1024 choices per position there are about a
billion possible codes for a catalog of 4.8 million products, so uniqueness
should be trivial. It is not. In our runs, roughly a quarter to a third of
products still share a code, depending on the method. The number of available
codes barely matters; what matters is how the code-assignment mechanism
distributes products over them.

The experiment compares three mechanisms under one budget on a new dataset,
TencentGR-1M (Tencent ads, all items embedded at 256 dimensions).
\textbf{RQ-VAE} learns the code system and the product map at the same time,
which is the standard recipe. \textbf{DRQ} learns the map first, freezes it, and
only then lays codes on top, with no feedback between the two phases. \textbf{Warmup}
splits the difference: map first for a few epochs, then the standard joint
learning, to test whether ``learning separately at all'' is what buys DRQ its
advantage, or only ``assigning codes by classic grouping instead of learning
them''. The results are a raw-numbers benchmark; deciding which mechanism wins
and why is the point of the study.

\section{Experiment design}

\subsection{Why re-implement}

No official implementation of DRQ is available: it was published by Shopee on a
proprietary corpus of 15M items, and neither the code nor the data accompany the
paper. The RQ-VAE side traces to TIGER (Google DeepMind), whose released code is
tied to its own training setup (Gumbel-Softmax selection, Sentence-T5 inputs on
Amazon review data). Neither is directly usable on our corpus. We therefore
re-implement both in a single script (\texttt{train\_drq\_vs\_rqvae.py}) that
exposes the three arms behind one flag (\texttt{--arm
rqvae\textbar warmup\textbar drq\textbar both\textbar all}), so every
difference in the outputs can only come from the mechanism itself, not from
budget or tuning. The paper's headline comparison is also on a different domain
(advertising embeddings versus e-commerce product text), so we benchmark the
claims on a new dataset rather than assume they transfer.

\subsection{Arms and shared budget}

All three arms see the identical configuration: latent dimension 64, $L = 3$
code levels, $K = 1024$ codes per level, 25 training epochs, batch 4096, Adam
with learning rate $5 \times 10^{-4}$, linear warmup over 2000 steps with cosine
decay, seed 42, and a 90/10 train/validation item split.

\begin{itemize}[nosep]
  \item \texttt{rqvae}: joint training of encoder, decoder, and learned
  codebooks with straight-through gradients (section~\ref{sec:rqvae}).
  Codebooks init truncated-normal; a dead-code revival pass re-seeds codes used
  below 10\% of mean usage. This is the baseline.
  \item \texttt{warmup}: 5 reconstruction-only epochs (the arm's encoder and
  decoder train with the quantizer bypassed), then the \emph{identical} joint
  training as \texttt{rqvae}. Control arm: it isolates ``clean latent before
  quantization'' from ``different assignment rule''.
  \item \texttt{drq}: two decoupled stages (section~\ref{sec:drq}). Stage 1 is
  the continuous VAE trained 25 epochs with $\lambda = 0$ (no contrastive term;
  the paper's DRQ-VAE variant). Stage 2 fits hierarchical $k$-means (FAISS, 50
  iterations, 3 restarts) on the frozen latents and assigns all items. The
  $\lambda > 0$ contrastive variant (DRQ-VAE+CL) is implemented but not yet run.
\end{itemize}

Every arm is then scored by the same diagnostics on all 4,783,154 items:
per-level codebook health (usage entropy, perplexity, Gini, effective codes
$\Keff$ under a swept kernel width), tuple-level injectivity (collision rate,
unique tuples, max collision bucket, effective cells), and validation
reconstruction, which keeps the comparison fair: a scheme that lowers collisions
by making worse codes has not really won.

\subsection{Dataset}

TencentGR-1M is the preliminary-round dataset of the Tencent Advertising
Algorithm Challenge 2025 (TAAC 2025, ``All-Modality Generative Recommendation''),
built from real, de-identified Tencent Ads logs and released publicly on
Hugging Face \emph{(org: TAAC2025)}. It provides roughly one million user
behavior sequences over a candidate pool of 4{,}783{,}154 ads, plus
organizer-supplied multi-modal item representations. We use the item side only:
a parquet table mapping each anonymous ad id to a 256-dimensional content
embedding, treated as a fixed input vector
$x \in \R^{256}$ that already fuses the ad's text and visual modalities. User
sequences, collaborative ids, and the remaining embedding tables (32 to 4096
dimensions) are not touched in this study.

The choice matters for the claims under test. DRQ's numbers come from a
proprietary Shopee corpus; TIGER's from Amazon review text. Advertising item
embeddings are a third, public, and much larger-than-paper regime per code space
(4.78M items against $1024^3$ nominal codes), which is exactly where
injectivity claims must be checked before they can be reused.

\section{Background: the two mechanisms}

\subsection{RQ-VAE: residual quantization inside a variational autoencoder}
\label{sec:rqvae}

Following SoundStream and TIGER, a quantizer $Q$ with $L$ levels maps a latent
$z \in \R^{d}$ to integer codes through a cascade of residual nearest-neighbor
searches. Level $\ell$ owns a codebook
$C^{(\ell)} = \{ e^{(\ell)}_i \}_{i=1}^{K} \subset \R^{d}$ of $K$ learned
centroids, and
\begin{align}
  r_0 &= z, \nonumber \\
  c_\ell &= \arg\min_{i} \; \norm{r_{\ell-1} - e^{(\ell)}_i}^2 ,
    \qquad
  r_\ell = r_{\ell-1} - e^{(\ell)}_{c_\ell}, \quad \ell = 1, \dots, L.
    \label{eq:rq}
\end{align}
The quantized latent is $\hat z = \sum_{\ell} e^{(\ell)}_{c_\ell}$ and the
semantic ID of an item is the tuple $(c_1, \dots, c_L)$. The tuple reads
coarse to fine because each level approximates the residual left by the levels
before it. Training minimizes reconstruction through a decoder $D$ plus the
vector-quantization objective
\begin{align}
  \mathcal{L}(x) \;=\; \norm{x - D(\hat z_{\mathrm{ste}})}^2
  \;+\; \sum_{\ell=1}^{L} \Bigl[
    \underbrace{\bigl\lVert \sg(r_{\ell-1}) - e^{(\ell)}_{c_\ell} \bigr\rVert^2}
      _{\text{codebook: pull centroid to its points}}
    \;+\; \beta \underbrace{\bigl\lVert r_{\ell-1}
      - \sg(e^{(\ell)}_{c_\ell}) \bigr\rVert^2}
      _{\text{commitment: pull point to its centroid}} \Bigr],
    \label{eq:loss}\\
  z_{\mathrm{ste}} &= z + \sg(\hat z - z), \label{eq:ste}
\end{align}
where $\sg(\cdot)$ stops gradients and $\beta$ (0.1 here) weights how hard the
encoder is pulled toward the frozen codebook. Equation~\eqref{eq:ste} is the
straight-through estimator: forward pass is quantized, backward pass copies
gradients at $\hat z$ through to $z$ as if $Q$ were the identity.

Two structural facts explain most of what follows. First, every reachable
$\hat z$ is a \emph{Minkowski sum} of the $L$ codebooks: the representable
latent set is the additive grid
$\bigl\{ \sum_\ell e^{(\ell)}_{i_\ell} \bigr\}$, which grows arithmetically as
the codebooks move under gradient descent. Second, a gradient reaches a centroid
only through samples for which it wins the $\arg\min$ in \eqref{eq:rq}; a code
that loses early gets a smaller share of updates, drifts further from the data,
and loses again. This is the failure mode the dead-code revival pass exists to
interrupt.

\subsection{Injectivity: why a billion slots still collide}

Let $N$ be the item count and $U$ the number of distinct full tuples among the
$N$ assignments. The benchmark's headline number is the collision rate
\begin{equation}
  \text{collision rate} \;=\; 1 - \frac{U}{N},
  \label{eq:collision}
\end{equation}
the fraction of items without a private ID. With $L = 3$, $K = 1024$, the
nominal tuple space is $K^L \approx 1.07 \times 10^9 \gg N = 4.78 \times 10^6$.
If items landed \emph{uniformly at random} over $M = K^L$ cells, the expected
unique fraction would follow the birthday relation
\begin{equation}
  \frac{U}{N} \;=\; \frac{1 - e^{-x}}{x}, \qquad x = \frac{N}{M},
  \label{eq:birthday}
\end{equation}
which at $x \approx 0.0045$ predicts $U/N \approx 0.998$: essentially
collision-free. Measured uniqueness around $0.72$ therefore cannot be a slot
shortage. It has two candidate sources. Skewed per-level usage caps how much of
the tuple space is reachable, since the reachable space is bounded by the
product of per-level \emph{effective} code counts. And joint concentration
wastes the rest: when an item's level-2 code mostly repeats what levels 0 and 1
already said, the cascade spends bits without buying separation.

DRQ's paper formalizes the per-level view with an expected overlap rate under a
Gaussian kernel of bandwidth $\sigma$, summed over used centroids,
\begin{equation}
  \Opi^{(\ell)}(\sigma) \;=\; \sum_{i}\sum_{j} \pi_i \pi_j\,
    \exp\!\Bigl(-\frac{\norm{c_i - c_j}^2}{4\sigma^2}\Bigr),
  \qquad
  \Keff^{(\ell)} \;=\; \frac{1}{\Opi^{(\ell)}},
  \label{eq:opiq}
\end{equation}
where $\pi_i$ is the usage fraction of code $i$. It splits exactly into a
\emph{distribution floor} $\sum_i \pi_i^2$ (usage skew) plus a \emph{geometry
cross-overlap} $\sum_{i \neq j} \pi_i \pi_j \exp(-\norm{c_i - c_j}^2 / 4\sigma^2)$
(centroids crowded onto the manifold). Balance fixes the first term only; the
second is what $k$-means-based assignment is argued to fix. We adopt
\eqref{eq:opiq} as the benchmark's geometry metric, reporting
$\Keff$ at a swept multiple of an auto-chosen $\sigma$ (0.1 $\times$ the median
nearest-neighbour distance between used centroids). In
\eqref{eq:opiq} the $c_i$ are the level's centroids; in \eqref{eq:rq} they are
the level-$\ell$ entries $e^{(\ell)}_i$. The metric is post-hoc: no diagnostic
parameter enters training or affects a single assigned ID.

\subsection{DRQ: decoupled residual quantization}
\label{sec:drq}

DRQ (arXiv:2606.01844) starts from the diagnosis above: in joint RQ-VAE training
the reconstruction loss wants $z$ to concentrate, the codebook objective wants
usage to spread, and both fight through the same gradients, so the geometry ends
up neither well spread nor well partitioned. It therefore splits training into
two stages that never share gradients.

\emph{Stage 1, continuous VAE.} Train $E, D$ with reconstruction only,
optionally plus a contrastive term:
\begin{equation}
  \mathcal{L}_{1}
  \;=\; \mathbb{E}_x \bigl\lVert x - D(E(x)) \bigr\rVert^2
  \;+\; \lambda\, \mathcal{L}_{\mathrm{contrastive}},
  \label{eq:drq-stage1}
\end{equation}
where $\mathcal{L}_{\mathrm{contrastive}}$ is a symmetric in-batch InfoNCE loss
on $z = E(x)$ between paired positive items (co-viewed or semantically similar;
temperature $\tau$), and $\lambda = 0$ gives the paper's DRQ-VAE variant while
$\lambda = 0.3$ gives DRQ-VAE+CL. No codebook, no argmin, no straight-through
pass: the latent geometry settles freely under the encoder's own objectives.

\emph{Stage 2, post-hoc hierarchical $k$-means.} Freeze the latents and build
the residual cascade of \eqref{eq:rq} level by level: level $\ell$ fits
$k$-means with $K$ centroids on the \emph{residual pool} left by levels
$1, \dots, \ell-1$ over all training items, then every item is assigned its
nearest centroid and passes its residual to the next level. Two properties
separate this from STE codebook training. Assignment quality is optimized
directly (each cell's centroids are the means of exactly the points assigned to
it), and every point votes on every centroid at every Lloyd iteration, unlike
\eqref{eq:loss} where only argmin-winners update. What stage 2 \emph{cannot} do
is reshape $z$: any pathology the VAE leaves in the latent is inherited.

\section{References}

\begin{enumerate}[nosep,leftmargin=*]
  \item DRQ: \emph{Rethinking Semantic ID Construction for Generative
        Recommendation} (arXiv:2606.01844). The arm reimplemented here; source
        of the overlap metric \eqref{eq:opiq}.
  \item TIGER: \emph{Recommender Systems with Generative Retrieval}
        (arXiv:2305.05065). The RQ-VAE semantic-ID recipe being benchmarked.
  \item SoundStream: \emph{An End-to-End Neural Audio Codec}
        (arXiv:2107.03312). Origin of residual quantization with the
        stop-gradient loss \eqref{eq:loss}.
  \item Lee et al., \emph{Autoregressive Image Generation using Residual
        Quantization} (CVPR 2022). The name ``RQ-VAE''.
  \item TencentGR-1M and the TAAC 2025 challenge:
        \texttt{huggingface.co/datasets/TAAC2025/TencentGR-1M}; competition
        paper arXiv:2604.04976.
  \item GRID handbook (arXiv:2507.22224). Independent evidence that RQ-VAE
        underperforms simpler SID construction; motivates the benchmark.
\end{enumerate}

\end{document}
